\chapter{2021年新高考I卷}

\section{单选题}

\begin{problem}
设集合 \(A=\{x\mid -2<x<4\}\)，\(B=\{2,3,4,5\}\)，则 \(A\cap B=\)（\quad）

\choices
    {\(\{2\}\)}
    {\(\{2,3\}\)}
    {\(\{3,4\}\)}
    {\(\{2,3,4\}\)}

\end{problem}

\begin{answer}
B
\end{answer}

\begin{solution}
由题设有 \(A\cap B=\{2,3\}\)，故选：B.
\end{solution}
\begin{problem}
已知 \(z=2-\i\)，则 \(z\left( \overline{z}+\i \right)=\)（\quad）

\choices
    {\(6-2\i\)}
    {\(4-2\i\)}
    {\(6+2\i\)}
    {\(4+2\i\)}

\end{problem}

\begin{answer}
C
\end{answer}

\begin{solution}
因为 \(z=2-\i\)，故 \(\overline{z}=2+\i\)，所以
\(z\left( \overline{z}+\i \right)=\left( 2-\i \right)\left( 2+2\i \right)=6+2\i\).

故选：C.
\end{solution}
\begin{problem}
已知圆锥的底面半径为 \(2\)，其侧面展开图为一个半圆，则该圆锥的母线长为（\quad）

\choices
    {\(2\)}
    {\(2\sqrt{2}\)}
    {\(4\)}
    {\(4\sqrt{2}\)}

\end{problem}

\begin{answer}
B
\end{answer}

\begin{solution}
设圆锥的母线长为 \(l\)，由于圆锥底面圆的周长等于扇形的弧长，则
\(\pi l=2\pi\times 2\)，
解得
\(l=2\sqrt{2}\).

故选：B.
\end{solution}
\begin{problem}
下列区间中，函数 \(f(x)=7\sin\left( x-\frac{\pi}{6}\right)\) 单调递增的区间是（\quad）

\choices
    {\(\left( 0,\frac{\pi}{2}\right)\)}
    {\(\left( \frac{\pi}{2},\pi\right)\)}
    {\(\left( \pi,\frac{3\pi}{2}\right)\)}
    {\(\left( \frac{3\pi}{2},2\pi\right)\)}

\end{problem}

\begin{answer}
A
\end{answer}

\begin{solution}
因为函数 \(y=\sin x\) 的单调递增区间为 \(\left( 2k\pi-\frac{\pi}{2},2k\pi+\frac{\pi}{2}\right)\)（\(k\in\Z\)），对于函数 \(f(x)=7\sin\left( x-\frac{\pi}{6}\right)\)，由 \(2k\pi-\frac{\pi}{2}<x-\frac{\pi}{6}<2k\pi+\frac{\pi}{2}\)（\(k\in\Z\)） 解得 \(2k\pi-\frac{\pi}{3}<x<2k\pi+\frac{2\pi}{3}\)（\(k\in\Z\)）.

取 \(k=0\)，可得函数 \(f(x)\) 的一个单调递增区间为 \(\left( -\frac{\pi}{3},\frac{2\pi}{3}\right)\).
则 \(\left( 0,\frac{\pi}{2}\right)\subset \left( -\frac{\pi}{3},\frac{2\pi}{3}\right)\)，\(\left( \frac{\pi}{2},\pi\right)\not\subset \left( -\frac{\pi}{3},\frac{2\pi}{3}\right)\).

取 \(k=1\)，可得函数 \(f(x)\) 的一个单调递增区间为 \(\left( \frac{5\pi}{3},\frac{8\pi}{3}\right)\).
因此 \(\left( \pi,\frac{3\pi}{2}\right)\not\subset \left( -\frac{\pi}{3},\frac{2\pi}{3}\right)\)，且 \(\left( \frac{3\pi}{2},2\pi\right)\not\subset \left( \frac{5\pi}{3},\frac{8\pi}{3}\right)\).

故选：A.
\end{solution}
\begin{problem}
已知 \(F_1,F_2\) 是椭圆 \(C:\frac{x^2}{9}+\frac{y^2}{4}=1\) 的两个焦点，点 \(M\) 在 \(C\) 上，则 \(MF_1\cdot MF_2\) 的最大值为（\quad）

\choices
    {\(13\)}
    {\(12\)}
    {\(9\)}
    {\(6\)}

\end{problem}

\begin{answer}
C
\end{answer}

\begin{solution}
由题，\(a^2=9\)，\(b^2=4\)，则
\(MF_1+MF_2=2a=6\).
所以
\(MF_1\cdot MF_2\le \left( \frac{MF_1+MF_2}{2}\right)^2=9\)
（当且仅当 \(MF_1=MF_2=3\) 时，等号成立）. 

故选：C.
\end{solution}
\begin{problem}
若 \(\tan\theta=-2\)，则 \(\dfrac{\sin\theta\left( 1+\sin2\theta \right)}{\sin\theta+\cos\theta}=\)（\quad）

\choices
    {\(-\dfrac{6}{5}\)}
    {\(-\dfrac{2}{5}\)}
    {\(\dfrac{2}{5}\)}
    {\(\dfrac{6}{5}\)}

\end{problem}

\begin{answer}
C
\end{answer}

\begin{solution}
将式子进行齐次化处理得：\(\frac{\sin\theta\left( 1+\sin2\theta \right)}{\sin\theta+\cos\theta}=\frac{\sin\theta\left( \sin\theta+\cos\theta+2\sin\theta\cos\theta \right)}{\sin\theta+\cos\theta}=\sin\theta\left( \sin\theta+\cos\theta \right)=\frac{\tan^2\theta+\tan\theta}{1+\tan^2\theta}\).

代入 \(\tan\theta=-2\)，得 \(\frac{\tan^2\theta+\tan\theta}{1+\tan^2\theta}=\frac{4-2}{1+4}=\frac{2}{5}\).

故选：C.
\end{solution}
\begin{problem}
若过点 \((a,b)\) 可以作曲线 \(y=\e^x\) 的两条切线，则（\quad）

\choices
    {\(\e^b<a\)}
    {\(\e^a<b\)}
    {\(0<a<\e^b\)}
    {\(0<b<\e^a\)}

\end{problem}

\begin{answer}
D
\end{answer}

\begin{solution}
在曲线 \(y=\e^x\) 上任取一点 \(P\left( t,\e^t \right)\)，对函数 \(y=\e^x\) 求导得 \(y'=\e^x\)，

所以，曲线 \(y=\e^x\) 在点 \(P\) 处的切线方程为
\(y-\e^t=\e^t\left( x-t \right)\)，
即
\(y=\e^t x+\left( 1-t \right)\e^t\).

由题意可知，点 \((a,b)\) 在直线
\(y=\e^t x+\left( 1-t \right)\e^t\)
上，可得
\(b=a\e^t+\left( 1-t \right)\e^t=\left( a+1-t \right)\e^t\).

令
\(f(t)=\left( a+1-t \right)\e^t\)，
则
\(f'(t)=\left( a-t \right)\e^t\).

当 \(t<a\) 时，\(f'(t)>0\)，此时函数 \(f(t)\) 单调递增；

当 \(t>a\) 时，\(f'(t)<0\)，此时函数 \(f(t)\) 单调递减；

所以，
\(f(t)_{\max}=f(a)=\e^a\).

由题意可知，直线 \(y=b\) 与曲线 \(y=f(t)\) 的图象有两个交点，则
\(b<f(t)_{\max}=\e^a\).

当 \(t<a+1\) 时，\(f(t)>0\)，当 \(t>a+1\) 时，\(f(t)<0\)，作出函数 \(f(t)\) 的图象如下图所示：

\bitmapfigure[max width=4.2cm,max totalheight=4.8cm]{img/2021/national_paper_1/q7_fig1.png}

由图可知，当
\(0<b<\e^a\)
时，直线 \(y=b\) 与曲线 \(y=f(t)\) 的图象有两个交点. 

故选：D.
\end{solution}
\begin{problem}
有 \(6\) 个相同的球，分别标有数字 \(1,2,3,4,5,6\)，从中有放回地随机取两次，每次取 \(1\) 个球，甲表示事件“第一次取出的球的数字是 \(1\)”，乙表示事件“第二次取出的球的数字是 \(2\)”，丙表示事件“两次取出的球的数字之和是 \(8\)”，丁表示事件“两次取出的球的数字之和是 \(7\)”，则（\quad）

\choices
    {甲与丙相互独立}
    {甲与丁相互独立}
    {乙与丙相互独立}
    {丙与丁相互独立}

\end{problem}

\begin{answer}
B
\end{answer}

\begin{solution}
根据独立事件概率关系逐一判断，\(P(\text{甲})=\frac{1}{6}\)，\(P(\text{乙})=\frac{1}{6}\)，\(P(\text{丙})=\frac{5}{36}\)，\(P(\text{丁})=\frac{6}{36}=\frac{1}{6}\).

又
\(P(\text{甲丙})=0\ne P(\text{甲})P(\text{丙})\)，
\(P(\text{甲丁})=\frac{1}{36}=P(\text{甲})P(\text{丁})\)，
\(P(\text{乙丙})=\frac{1}{36}\ne P(\text{乙})P(\text{丙})\)，
\(P(\text{丙丁})=0\ne P(\text{丁})P(\text{丙})\).

故选：B.
\end{solution}
\section{多选题}

\begin{problem}
有一组样本数据 \(x_1,x_2,\dots,x_n\)，由这组数据得到新样本数据 \(y_1,y_2,\dots,y_n\)，其中 \(y_i=x_i+c\)（\(i=1,2,\dots,n\)），\(c\) 为非零常数，则（\quad）

\choices
    {两组样本数据的样本平均数相同}
    {两组样本数据的样本中位数相同}
    {两组样本数据的样本标准差相同}
    {两组样数据的样本极差相同}

\end{problem}

\begin{answer}
CD
\end{answer}

\begin{solution}
A：\(E(y)=E(x+c)=E(x)+c\)，且 \(c\ne 0\)，故平均数不相同，错误；

B：设第一组样本数据的中位数为 \(m\)，则第二组样本数据的中位数为 \(m+c\). 又 \(c\ne 0\)，所以 \(m+c\ne m\)，故两组样本数据的样本中位数不相同，错误；

C：\(D(y)=D(x)+D(c)=D(x)\)，故方差相同，正确；

D：由极差的定义知：若第一组的极差为 \(x_{\max}-x_{\min}\)，则第二组的极差为 \(y_{\max}-y_{\min}=\left( x_{\max}+c \right)-\left( x_{\min}+c \right)=x_{\max}-x_{\min}\)，故极差相同，正确. 

故选：CD.
\end{solution}
\begin{problem}
已知 \(O\) 为坐标原点，点 \(P_1=\left( \cos\alpha,\sin\alpha \right)\)，\(P_2=\left( \cos\beta,-\sin\beta \right)\)，\(P_3=\left( \cos\left( \alpha+\beta \right),\sin\left( \alpha+\beta \right) \right)\)，\(A(1,0)\)，则（\quad）

\choices
    {\(|OP_1|=|OP_2|\)}
    {\(|AP_1|=|AP_2|\)}
    {\(\overrightarrow{OA}\cdot\overrightarrow{OP_3}=\overrightarrow{OP_1}\cdot\overrightarrow{OP_2}\)}
    {\(\overrightarrow{OA}\cdot\overrightarrow{OP_1}=\overrightarrow{OP_2}\cdot\overrightarrow{OP_3}\)}

\end{problem}

\begin{answer}
AC
\end{answer}

\begin{solution}
A：\(\overrightarrow{OP_1}=\left( \cos\alpha,\sin\alpha \right)\)，\(\overrightarrow{OP_2}=\left( \cos\beta,-\sin\beta \right)\)，所以 \(|OP_1|=\sqrt{\cos^2\alpha+\sin^2\alpha}=1\)， \(|OP_2|=\sqrt{\cos^2\beta+\left( -\sin\beta \right)^2}=1\)， 故 \(|OP_1|=|OP_2|\)，正确；

B：\(\overrightarrow{AP_1}=\left( \cos\alpha-1,\sin\alpha \right)\)，\(\overrightarrow{AP_2}=\left( \cos\beta-1,-\sin\beta \right)\)，所以 \(|AP_1|=\sqrt{\left( \cos\alpha-1 \right)^2+\sin^2\alpha}=2\left| \sin\frac{\alpha}{2}\right|\)， \(|AP_2|=\sqrt{\left( \cos\beta-1 \right)^2+\sin^2\beta}=2\left| \sin\frac{\beta}{2}\right|\)， 故 \(|AP_1|\)、\(|AP_2|\) 不一定相等，错误；

C：由题意得 \(\overrightarrow{OA}\cdot\overrightarrow{OP_3}=1\times \cos\left( \alpha+\beta \right)+0\times \sin\left( \alpha+\beta \right)=\cos\left( \alpha+\beta \right)\)， \(\overrightarrow{OP_1}\cdot\overrightarrow{OP_2}=\cos\alpha\cos\beta+\sin\alpha\left( -\sin\beta \right)=\cos\left( \alpha+\beta \right)\)， 正确；

D：由题意得 \(\overrightarrow{OA}\cdot\overrightarrow{OP_1}=\cos\alpha\)， 而 \(\overrightarrow{OP_2}\cdot\overrightarrow{OP_3}=\cos\beta\cos\left( \alpha+\beta \right)+\left( -\sin\beta \right)\sin\left( \alpha+\beta \right)=\cos\left( \alpha+2\beta \right)\)， 一般不相等，错误. 

故选：AC.
\end{solution}
\begin{problem}
已知点 \(P\) 在圆 \(\left( x-5 \right)^2+\left( y-5 \right)^2=16\) 上，点 \(A(4,0)\)、\(B(0,2)\)，则（\quad）

\choices
    {点 \(P\) 到直线 \(AB\) 的距离小于 \(10\)}
    {点 \(P\) 到直线 \(AB\) 的距离大于 \(2\)}
    {当 \(\angle PBA\) 最小时，\(PB=3\sqrt{2}\)}
    {当 \(\angle PBA\) 最大时，\(PB=3\sqrt{2}\)}

\end{problem}

\begin{answer}
ACD
\end{answer}

\begin{solution}
圆 \(\left( x-5 \right)^2+\left( y-5 \right)^2=16\) 的圆心为 \(M(5,5)\)，半径为 \(4\). 

直线 \(AB\) 的方程为
\(\frac{x}{4}+\frac{y}{2}=1\)，
即
\(x+2y-4=0\).

圆心 \(M\) 到直线 \(AB\) 的距离为 \(\frac{|5+2\times 5-4|}{\sqrt{1^2+2^2}}=\frac{11}{\sqrt{5}}=\frac{11\sqrt{5}}{5}>4\).

所以，点 \(P\) 到直线 \(AB\) 的距离的最小值为
\(\frac{11\sqrt{5}}{5}-4<2\)，
最大值为
\(\frac{11\sqrt{5}}{5}+4<10\)，
A 选项正确，B 选项错误；

如下图所示：

\bitmapfigure[max width=5.2cm,max totalheight=4.8cm]{img/2021/national_paper_1/q11_fig1.png}

当 \(\angle PBA\) 最大或最小时，\(PB\) 与圆 \(M\) 相切，连接 \(MP\)、\(BM\)，可知 \(PM\perp PB\)，
\(BM=\sqrt{\left( 0-5 \right)^2+\left( 2-5 \right)^2}=\sqrt{34}\)，
\(MP=4\)，
由勾股定理可得
\(PB=\sqrt{BM^2-MP^2}=3\sqrt{2}\).

故 C、D 选项正确. 

故选：ACD.
\end{solution}
\begin{problem}
在正三棱柱 \(ABC-A_1B_1C_1\) 中，\(AB=AA_1=1\)，点 \(P\) 满足 \(\overrightarrow{BP}=\lambda\overrightarrow{BC}+\mu\overrightarrow{BB_1}\)，其中 \(\lambda,\mu\in[0,1]\).

\bitmapfigure[max width=3.4cm,max totalheight=4.8cm]{img/2021/national_paper_1/q12_fig1.png}

则（\quad）

\choices
    {当 \(\lambda=1\) 时，\(\triangle AB_1P\) 的周长为定值}
    {当 \(\mu=1\) 时，三棱锥 \(P-A_1BC\) 的体积为定值}
    {当 \(\lambda=\frac{1}{2}\) 时，有且仅有一个点 \(P\)，使得 \(A_1P\perp BP\)}
    {当 \(\mu=\frac{1}{2}\) 时，有且仅有一个点 \(P\)，使得 \(A_1B\perp\) 平面 \(AB_1P\)}

\end{problem}

\begin{answer}
BD
\end{answer}

\begin{solution}
因为 \(\overrightarrow{BP}=\lambda\overrightarrow{BC}+\mu\overrightarrow{BB_1}\)，且 \(\lambda,\mu\in[0,1]\)，所以点 \(P\) 由点 \(B\) 沿 \(\overrightarrow{BC}\) 与 \(\overrightarrow{BB_1}\) 的方向在不超过对应边长的范围内平移得到，故点 \(P\) 在矩形 \(BCC_1B_1\) 内部（含边界）. 

对于 A，当 \(\lambda=1\) 时，\(\overrightarrow{BP}=\overrightarrow{BC}+\mu\overrightarrow{BB_1}=\overrightarrow{BC}+\mu\overrightarrow{CC_1}\)，即此时 \(P\in \overline{CC_1}\)，\(\triangle AB_1P\) 周长不是定值，故 A 错误；

对于 B，当 \(\mu=1\) 时，\(\overrightarrow{BP}=\lambda\overrightarrow{BC}+\overrightarrow{BB_1}=\overrightarrow{BB_1}+\lambda\overrightarrow{B_1C_1}\)，故此时 \(P\) 点轨迹为线段 \(\overline{B_1C_1}\)，而 \(B_1C_1 // BC\)，且 \(B_1C_1 // \) 平面 \(A_1BC\)，则 \(P\) 到平面 \(A_1BC\) 的距离为定值，所以其体积为定值，故 B 正确；

对于 C，当 \(\lambda=\frac{1}{2}\) 时，\(\overrightarrow{BP}=\frac{1}{2}\overrightarrow{BC}+\mu\overrightarrow{BB_1}\). 取 \(BC\)、\(B_1C_1\) 中点分别为 \(Q\)、\(H\)，则 \(\overrightarrow{BP}=\overrightarrow{BQ}+\mu\overrightarrow{QH}\)，所以 \(P\) 点轨迹为线段 \(\overline{QH}\). 不妨建立空间直角坐标系如图，则 \(A_1\left( \frac{\sqrt{3}}{2},0,1 \right)\)，\(P(0,0,\mu)\)，\(B\left( 0,\frac{1}{2},0 \right)\).
于是 \(\overrightarrow{A_1P}=\left( -\frac{\sqrt{3}}{2},0,\mu-1 \right)\)，\(\overrightarrow{BP}=\left( 0,-\frac{1}{2},\mu \right)\).
由 \(A_1P\perp BP\) 得 \(\mu\left( \mu-1 \right)=0\)，所以 \(\mu=0\) 或 \(\mu=1\)，故 \(Q\)、\(H\) 均满足，故 C 错误；

对于 D，当 \(\mu=\frac{1}{2}\) 时，\(\overrightarrow{BP}=\lambda\overrightarrow{BC}+\frac{1}{2}\overrightarrow{BB_1}\). 取 \(BB_1\)、\(CC_1\) 中点为 \(M\)、\(N\)，则 \(\overrightarrow{BP}=\overrightarrow{BM}+\lambda\overrightarrow{MN}\)，所以 \(P\) 点轨迹为线段 \(\overline{MN}\). 设 \(P\left( 0,y_0,\frac{1}{2} \right)\)，\(A\left( \frac{\sqrt{3}}{2},0,0 \right)\). 则 \(\overrightarrow{AP}=\left( -\frac{\sqrt{3}}{2},y_0,\frac{1}{2} \right)\)，\(\overrightarrow{A_1B}=\left( -\frac{\sqrt{3}}{2},\frac{1}{2},-1 \right)\).
又 \(\overrightarrow{A_1B}\perp \overrightarrow{AB_1}\)，故只需 \(\overrightarrow{A_1B}\perp \overrightarrow{AP}\). 

由
\(\overrightarrow{A_1B}\cdot \overrightarrow{AP}=0\)
得
\(\frac{3}{4}+\frac{1}{2}y_0-\frac{1}{2}=0\)，
即
\(y_0=-\frac{1}{2}\).
此时 \(P\) 与 \(N\) 重合，故满足条件的点 \(P\) 有且仅有一个，故 D 正确. 

故选：BD.
\end{solution}
\section{填空题}

\begin{problem}
已知函数 \(f(x)=x^3\left( a\cdot 2^x-2^{-x} \right)\) 是偶函数，则 \(a=\fillinblank{}\).

\end{problem}

\begin{answer}
1
\end{answer}

\begin{solution}
因为
\(f(x)=x^3\left( a\cdot 2^x-2^{-x} \right)\)，
故
\(f(-x)=-x^3\left( a\cdot 2^{-x}-2^x \right)\).

因为 \(f(x)\) 为偶函数，故 \(f(-x)=f(x)\)，

于是
\(x^3\left( a\cdot 2^x-2^{-x} \right)=-x^3\left( a\cdot 2^{-x}-2^x \right)\)，
整理得
\(\left( a-1 \right)\left( 2^x+2^{-x} \right)=0\).

故 \(a=1\).
\end{solution}
\begin{problem}
已知 \(O\) 为坐标原点，抛物线 \(C:y^2=2px\)（\(p>0\)）的焦点为 \(F\)，\(P\) 为 \(C\) 上一点，\(PF\) 与 \(x\) 轴垂直，\(Q\) 为 \(x\) 轴上一点，且 \(PQ\perp OP\)，若 \(FQ=6\)，则 \(C\) 的准线方程为 \fillinblank{}.

\end{problem}

\begin{answer}
\(x=-\frac{3}{2}\)
\end{answer}

\begin{solution}
不妨设
\(P\left( \frac{p}{2},p \right)\)，\(Q\left( 6+\frac{p}{2},0 \right)\)，
则
\(\overrightarrow{PQ}=(6,-p)\).

因为 \(\overrightarrow{PQ}\perp \overrightarrow{OP}\)，所以
\(\left( \frac{p}{2},p \right)\cdot (6,-p)=0\).
即
\(\frac{p}{2}\times 6-p^2=0\).

又 \(p>0\)，解得
\(p=3\).

所以抛物线 \(C\) 的准线方程为
\(x=-\frac{p}{2}=-\frac{3}{2}\).
\end{solution}
\begin{problem}
函数 \(f(x)=|2x-1|-2\ln x\) 的最小值为 \fillinblank{}.

\end{problem}

\begin{answer}
1
\end{answer}

\begin{solution}
由题设知 \(f(x)=|2x-1|-2\ln x\) 的定义域为 \((0,+\infty)\). 

当 \(0<x\le \frac{1}{2}\) 时，
\(f(x)=1-2x-2\ln x\)，
此时 \(f(x)\) 单调递减；

当 \(\frac{1}{2}<x\le 1\) 时，
\(f(x)=2x-1-2\ln x\)，
有
\(f'(x)=2-\frac{2}{x}\le 0\)，
此时 \(f(x)\) 单调递减；

当 \(x>1\) 时，
\(f(x)=2x-1-2\ln x\)，
有
\(f'(x)=2-\frac{2}{x}>0\)，
此时 \(f(x)\) 单调递增. 

又 \(f(x)\) 在各分段的界点处连续，

综上有：\(0<x\le 1\) 时，\(f(x)\) 单调递减；\(x>1\) 时，\(f(x)\) 单调递增，

故
\(f(x)\ge f(1)=1\).
\end{solution}
\begin{problem}
某校学生在研究民间剪纸艺术时，发现剪纸时经常会沿纸的某条对称轴把纸对折，规格为 \(20\text{dm}\times 12\text{dm}\) 的长方形纸，对折 \(1\) 次共可以得到 \(10\text{dm}\times 12\text{dm}\)、\(20\text{dm}\times 6\text{dm}\) 两种规格的图形，它们的面积之和 \(S_1=240\text{dm}^2\)；对折 \(2\) 次共可以得到 \(5\text{dm}\times 12\text{dm}\)、\(10\text{dm}\times 6\text{dm}\)、\(20\text{dm}\times 3\text{dm}\) 三种规格的图形，它们的面积之和 \(S_2=180\text{dm}^2\). 以此类推.

\begin{enumerate}
\item 对折 \(4\) 次共可以得到不同规格图形的种数为 \fillinblank{}；
\item 如果对折 \(n\) 次，那么 \(\sum_{k=1}^n S_k=\fillinblank{}\text{dm}^2\).
\end{enumerate}

\end{problem}

\begin{answer}
5；\(720-\dfrac{15(n+3)}{2^{n-4}}\)
\end{answer}

\begin{solution}
（1）对折 \(4\) 次可得到如下规格：\(\frac{5}{4}\text{dm}\times 12\text{dm}\)，\(\frac{5}{2}\text{dm}\times 6\text{dm}\)，\(5\text{dm}\times 3\text{dm}\)，\(10\text{dm}\times \frac{3}{2}\text{dm}\)，\(20\text{dm}\times \frac{3}{4}\text{dm}\)，共 \(5\) 种；

（2）由题意可得 \(S_1=2\times 120,\quad S_2=3\times 60,\quad S_3=4\times 30,\quad S_4=5\times 15,\ \dots\)， 故 \(S_n=\frac{120(n+1)}{2^{n-1}}\).

设 \(S=\frac{120\times 2}{2^0}+\frac{120\times 3}{2^1}+\frac{120\times 4}{2^2}+\cdots+\frac{120(n+1)}{2^{n-1}}\)， 则 \(\frac{1}{2}S=\frac{120\times 2}{2^1}+\frac{120\times 3}{2^2}+\cdots+\frac{120n}{2^{n-1}}+\frac{120(n+1)}{2^n}\).

两式作差得 \(\frac{1}{2}S=240+120\left( \frac{1}{2}+\frac{1}{2^2}+\cdots+\frac{1}{2^{n-1}} \right)-\frac{120(n+1)}{2^n}\).

而 \(\frac{1}{2}+\frac{1}{2^2}+\cdots+\frac{1}{2^{n-1}}=1-\frac{1}{2^{n-1}}\)， 所以 \(\frac{1}{2}S=240+120\left( 1-\frac{1}{2^{n-1}} \right)-\frac{120(n+1)}{2^n}=360-\frac{120}{2^{n-1}}-\frac{120(n+1)}{2^n}\).

即 \(\frac{1}{2}S=360-\frac{120(n+3)}{2^n}\)， 从而 \(S=720-\frac{240(n+3)}{2^n}=720-\frac{15(n+3)}{2^{n-4}}\).
\end{solution}
\section{解答题}

\begin{problem}
已知数列 \(\{a_n\}\) 满足 \(a_1=1\)，当 \(n\) 为奇数时，\(a_{n+1}=a_n+1\)；当 \(n\) 为偶数时，\(a_{n+1}=a_n+2\).

\begin{enumerate}
\item 记 \(b_n=a_{2n}\)，写出 \(b_1\)、\(b_2\)，并求数列 \(\{b_n\}\) 的通项公式；
\item 求 \(\{a_n\}\) 的前 \(20\) 项和.
\end{enumerate}

\end{problem}

\begin{solution}
（1）由题设可得
\(b_1=a_2=a_1+1=2\)，\(b_2=a_4=a_3+1=a_2+2+1=5\).

又
\(a_{2k+2}=a_{2k+1}+1\)，\(a_{2k+1}=a_{2k}+2\)，
故
\(a_{2k+2}=a_{2k}+3\)，
即
\(b_{n+1}=b_n+3\).

所以 \(\{b_n\}\) 为等差数列，故
\(b_n=2+\left( n-1 \right)\times 3=3n-1\).

（2）设 \(\{a_n\}\) 的前 \(20\) 项和为 \(S_{20}\)，则
\(S_{20}=a_1+a_2+a_3+\cdots+a_{20}\).

因为
\(a_1=a_2-1,\quad a_3=a_4-1,\quad \dots,\quad a_{19}=a_{20}-1\)，
所以 \(S_{20}=2\left( a_2+a_4+\cdots+a_{20} \right)-10=2\left( b_1+b_2+\cdots+b_{10} \right)-10\).

由（1）知
\(b_n=3n-1\)，
故
\(S_{20}=2\times \frac{\left( 2+29 \right)\times 10}{2}-10=300\).
\end{solution}
\begin{problem}
某学校组织“一带一路”知识竞赛，有 \(A\)、\(B\) 两类问题，每位参加比赛的同学先在两类问题中选择一类并从中随机抽取一个问题回答，若回答错误则该同学比赛结束；若回答正确则从另一类问题中再随机抽取一个问题回答，无论回答正确与否，该同学比赛结束. \(A\) 类问题中的每个问题回答正确得 \(20\) 分，否则得 \(0\) 分；\(B\) 类问题中的每个问题回答正确得 \(80\) 分，否则得 \(0\) 分. 已知小明能正确回答 \(A\) 类问题的概率为 \(0.8\)，能正确回答 \(B\) 类问题的概率为 \(0.6\)，且能正确回答问题的概率与回答次序无关.

\begin{enumerate}
\item 若小明先回答 \(A\) 类问题，记 \(X\) 为小明的累计得分，求 \(X\) 的分布列；
\item 为使累计得分的期望最大，小明应选择先回答哪类问题？并说明理由.
\end{enumerate}

\end{problem}

\begin{solution}
（1）由题可知，\(X\) 的所有可能取值为 \(0\)、\(20\)、\(100\). 

\(P(X=0)=1-0.8=0.2\)，
\(P(X=20)=0.8\left( 1-0.6 \right)=0.32\)，
\(P(X=100)=0.8\times 0.6=0.48\).

所以 \(X\) 的分布列为

\begin{center}
\begin{tabular}{c|ccc}
\hline
\(X\) & \(0\) & \(20\) & \(100\) \\
\hline
\(P\) & \(0.2\) & \(0.32\) & \(0.48\) \\
\hline
\end{tabular}
\end{center}

（2）由（1）知
\(E(X)=0\times 0.2+20\times 0.32+100\times 0.48=54.4\).

若小明先回答 \(B\) 类问题，记 \(Y\) 为小明的累计得分，则 \(Y\) 的所有可能取值为 \(0\)、\(80\)、\(100\). 

\(P(Y=0)=1-0.6=0.4\)，
\(P(Y=80)=0.6\left( 1-0.8 \right)=0.12\)，
\(P(Y=100)=0.8\times 0.6=0.48\).

所以
\(E(Y)=0\times 0.4+80\times 0.12+100\times 0.48=57.6\).

因为
\(54.4<57.6\)，
所以小明应选择先回答 \(B\) 类问题.
\end{solution}
\begin{problem}
记 \(\triangle ABC\) 的内角 \(A\)、\(B\)、\(C\) 的对边分别为 \(a\)、\(b\)、\(c\). 已知 \(b^2=ac\)，点 \(D\) 在边 \(AC\) 上，且 \(BD\sin\angle ABC=a\sin C\). 

\begin{enumerate}
\item 证明：\(BD=b\)；
\item 若 \(AD=2DC\)，求 \(\cos\angle ABC\).
\end{enumerate}

\bitmapfigure[max width=5.0cm,max totalheight=4.8cm]{img/2021/national_paper_1/q19_fig1.png}

\end{problem}

\begin{solution}
（1）由题设，
\(BD=\frac{a\sin C}{\sin\angle ABC}\).
由正弦定理知
\(\frac{c}{\sin C}=\frac{b}{\sin\angle ABC}\)，
即
\(\frac{\sin C}{\sin\angle ABC}=\frac{c}{b}\).
所以
\(BD=\frac{ac}{b}\).

又 \(b^2=ac\)，故
\(BD=b\).

（2）由题意知
\(BD=b\)，\(AD=\frac{2b}{3}\)，\(DC=\frac{b}{3}\).

在 \(\triangle ADB\) 中，由余弦定理得 \(\cos\angle ADB=\frac{AD^2+BD^2-AB^2}{2AD\cdot BD}=\frac{\frac{4b^2}{9}+b^2-c^2}{2\times \frac{2b}{3}\times b}=\frac{13b^2-9c^2}{12b^2}\).

在 \(\triangle CDB\) 中，由余弦定理得 \(\cos\angle CDB=\frac{DC^2+BD^2-BC^2}{2DC\cdot BD}=\frac{\frac{b^2}{9}+b^2-a^2}{2\times \frac{b}{3}\times b}=\frac{10b^2-9a^2}{6b^2}\).

因为
\(\angle ADB=\pi-\angle CDB\)，
所以
\(\cos\angle ADB=-\cos\angle CDB\).
即
\(\frac{13b^2-9c^2}{12b^2}=-\frac{10b^2-9a^2}{6b^2}\)，
整理得
\(2a^2+\frac{b^4}{a^2}=\frac{11b^2}{3}\).

设 \(x=\frac{a^2}{b^2}\)，则
\(2x^2-\frac{11}{3}x+1=0\).
解得
\(x=\frac{1}{3}\)或\(x=\frac{3}{2}\).

由余弦定理，
\(\cos\angle ABC=\frac{a^2+c^2-b^2}{2ac}\).
又 \(b^2=ac\)，所以
\(\cos\angle ABC=\frac{a^2+\frac{b^4}{a^2}-b^2}{2b^2}\).

当 \(\frac{a^2}{b^2}=\frac{1}{3}\) 时，
\(\cos\angle ABC=\frac{7}{6}>1\)，
不合题意；

当 \(\frac{a^2}{b^2}=\frac{3}{2}\) 时，
\(\cos\angle ABC=\frac{7}{12}\).
\end{solution}
\begin{problem}
如图，在三棱锥 \(A-BCD\) 中，平面 \(ABD\perp\) 平面 \(BCD\)，\(AB=AD\)，\(O\) 为 \(BD\) 的中点.

\begin{enumerate}
\item 证明：\(OA\perp CD\)；
\item 若 \(\triangle OCD\) 是边长为 \(1\) 的等边三角形，点 \(E\) 在棱 \(AD\) 上，\(DE=2EA\)，且二面角 \(E-BC-D\) 的大小为 \(45^\circ\)，求三棱锥 \(A-BCD\) 的体积.
\end{enumerate}

\bitmapfigure[max width=4.6cm,max totalheight=4.8cm]{img/2021/national_paper_1/q20_fig1.png}

\end{problem}

\begin{solution}
（1）因为 \(AB=AD\)，\(O\) 为 \(BD\) 中点，所以
\(AO\perp BD\).

因为平面 \(ABD\cap\) 平面 \(BCD=BD\)，平面 \(ABD\perp\) 平面 \(BCD\)，且 \(AO\subset\) 平面 \(ABD\)，

因此
\(AO\perp \text{平面 }BCD\).

因为 \(CD\subset\) 平面 \(BCD\)，所以
\(AO\perp CD\).

（2）作 \(EF\perp BD\) 于 \(F\)，作 \(FM\perp BC\) 于 \(M\)，连 \(EM\)、\(FM\). 

\bitmapfigure[max width=5.5cm,max totalheight=4.8cm]{img/2021/national_paper_1/q20_fig2.png}

因为 \(AO\perp\) 平面 \(BCD\)，所以
\(AO\perp BD\)，\(AO\perp CD\).

又 \(EF // AO\)，故
\(EF\perp BD\)，\(EF\perp CD\).
且 \(BD\cap CD=D\)，因此
\(EF\perp \text{平面 }BCD\).
所以
\(EF\perp BC\).

因为 \(FM\perp BC\)，且 \(FM\cap EF=F\)，所以
\(BC\perp \text{平面 }EFM\).
故
\(BC\perp MF\).
则 \(\angle EMF\) 为二面角 \(E-BC-D\) 的平面角，所以
\(\angle EMF=\frac{\pi}{4}\).

因为 \(BO=OD\)，且 \(\triangle OCD\) 是边长为 \(1\) 的等边三角形，所以
\(OD=OC=CD=1\)，\(BO=1\)，\(BD=2\).
从而 \(\triangle BCD\) 是直角三角形，且
\(BC=\sqrt{BD^2-CD^2}=\sqrt{4-1}=\sqrt{3}\).

又 \(DE=2EA\)，所以
\(\frac{DE}{DA}=\frac{2}{3}\).
由 \(EF // AO\)，得
\(\frac{EF}{AO}=\frac{DE}{DA}=\frac{2}{3}\).

在直角三角形 \(BCD\) 中，因 \(FM // CD\)，有
\(\frac{FM}{CD}=\frac{BF}{BD}\).
又
\(DF=\frac{2}{3}DO=\frac{2}{3}\)，
所以
\(BF=BD-DF=2-\frac{2}{3}=\frac{4}{3}\).
从而
\(FM=\frac{BF}{BD}\cdot CD=\frac{\frac{4}{3}}{2}\cdot 1=\frac{2}{3}\).

因为 \(\angle EMF=45^\circ\)，且 \(EF\perp FM\)，所以
\(EF=FM=\frac{2}{3}\).
于是
\(AO=\frac{3}{2}EF=1\).

底面 \(\triangle BCD\) 的面积为
\(S_{\triangle BCD}=\frac{1}{2}\times BC\times CD=\frac{\sqrt{3}}{2}\).

因此三棱锥 \(A-BCD\) 的体积 \(V=\frac{1}{3}\cdot AO\cdot S_{\triangle BCD}=\frac{1}{3}\times 1\times \frac{\sqrt{3}}{2}=\frac{\sqrt{3}}{6}\).
\end{solution}
\begin{problem}
在平面直角坐标系 \(xOy\) 中，已知点 \(F_1(-\sqrt{17},0)\)、\(F_2(\sqrt{17},0)\)，满足 \(MF_1-MF_2=2\) 的点 \(M\) 的轨迹为 \(C\). 

\begin{enumerate}
\item 求 \(C\) 的方程；
\item 设点 \(T\) 在直线 \(x=\frac{1}{2}\) 上，过 \(T\) 的两条直线分别交 \(C\) 于 \(A\)、\(B\) 两点和 \(P\)、\(Q\) 两点，且 \(TA\cdot TB=TP\cdot TQ\)，求直线 \(AB\) 的斜率与直线 \(PQ\) 的斜率之和.
\end{enumerate}

\end{problem}

\begin{solution}
（1）因为
\(MF_1-MF_2=2<2\sqrt{17}=F_1F_2\)，
所以，轨迹 \(C\) 是以点 \(F_1\)、\(F_2\) 为左、右焦点的双曲线的右支. 

设轨迹 \(C\) 的方程为
\(\frac{x^2}{a^2}-\frac{y^2}{b^2}=1\)（\(a>0\)，\(b>0\)），
则
\(2a=2\)，
可得
\(a=1\).
又
\(c=\sqrt{17}\)，
所以
\(b^2=c^2-a^2=17-1=16\).

故轨迹 \(C\) 的方程为
\(x^2-\frac{y^2}{16}=1\)（\(x\ge 1\)）.

（2）设点
\(T\left( \frac{1}{2},t \right)\).
若过点 \(T\) 的直线斜率不存在，此时该直线与曲线 \(C\) 无公共点，不妨设直线 \(AB\) 的斜率为 \(k_1\)，则其方程为
\(y-t=k_1\left( x-\frac{1}{2} \right)\)，
即
\(y=k_1x+t-\frac{k_1}{2}\).

联立 \(\begin{cases}
y=k_1x+t-\frac{k_1}{2},\\
16x^2-y^2=16,
\end{cases}\)，消去 \(y\) 并整理可得
\(\left( k_1^2-16 \right)x^2+k_1\left( 2t-k_1 \right)x+\left( t-\frac{k_1}{2} \right)^2+16=0\).

设点 \(A(x_1,y_1)\)、\(B(x_2,y_2)\)，则由韦达定理得
\(x_1+x_2=-\frac{k_1\left( 2t-k_1 \right)}{k_1^2-16}\)，
\(x_1x_2=\frac{\left( t-\frac{k_1}{2} \right)^2+16}{k_1^2-16}\).

所以 \(TA\cdot TB=\left( 1+k_1^2 \right)\left( x_1-\frac{1}{2} \right)\left( x_2-\frac{1}{2} \right)=\left( 1+k_1^2 \right)\left( x_1x_2-\frac{x_1+x_2}{2}+\frac{1}{4} \right)=\frac{\left( t^2+12 \right)\left( 1+k_1^2 \right)}{k_1^2-16}\).

设直线 \(PQ\) 的斜率为 \(k_2\)，则其方程为
\(y-t=k_2\left( x-\frac{1}{2} \right)\)，
即
\(y=k_2x+t-\frac{k_2}{2}\).

联立 \(\begin{cases}
y=k_2x+t-\frac{k_2}{2},\\
16x^2-y^2=16,
\end{cases}\)，消去 \(y\) 并整理可得
\(\left( k_2^2-16 \right)x^2+k_2\left( 2t-k_2 \right)x+\left( t-\frac{k_2}{2} \right)^2+16=0\).

设点 \(P(x_3,y_3)\)、\(Q(x_4,y_4)\)，则由韦达定理得
\(x_3+x_4=-\frac{k_2\left( 2t-k_2 \right)}{k_2^2-16}\)，
\(x_3x_4=\frac{\left( t-\frac{k_2}{2} \right)^2+16}{k_2^2-16}\).

所以 \(TP\cdot TQ=\left( 1+k_2^2 \right)\left( x_3-\frac{1}{2} \right)\left( x_4-\frac{1}{2} \right)=\left( 1+k_2^2 \right)\left( x_3x_4-\frac{x_3+x_4}{2}+\frac{1}{4} \right)=\frac{\left( t^2+12 \right)\left( 1+k_2^2 \right)}{k_2^2-16}\).

因为 \(TA\cdot TB=TP\cdot TQ\)，即
\(\frac{\left( t^2+12 \right)\left( 1+k_1^2 \right)}{k_1^2-16}=\frac{\left( t^2+12 \right)\left( 1+k_2^2 \right)}{k_2^2-16}\)，整理可得
\(\left( k_1-k_2 \right)\left( k_1+k_2 \right)=0\).

由于过点 \(T\) 且斜率为给定值的直线只有一条，而直线 \(AB\) 与直线 \(PQ\) 是两条直线，所以 \(k_1\ne k_2\)，故
\(k_1+k_2=0\).

因此，直线 \(AB\) 的斜率与直线 \(PQ\) 的斜率之和为
\(0\).
\end{solution}
\begin{problem}
已知函数 \(f(x)=x\left( 1-\ln x \right)\). 

\begin{enumerate}
\item 讨论 \(f(x)\) 的单调性；
\item 设 \(a\)、\(b\) 为两个不相等的正数，且 \(b\ln a-a\ln b=a-b\)，证明 \(2<\frac{1}{a}+\frac{1}{b}<\e\). 
\end{enumerate}

\end{problem}

\begin{solution}
（1）函数的定义域为 \((0,+\infty)\)，

又
\(f'(x)=1-\ln x-1=-\ln x\).

当 \(x\in(0,1)\) 时，\(f'(x)>0\)；当 \(x\in(1,+\infty)\) 时，\(f'(x)<0\)，

故 \(f(x)\) 的递增区间为 \((0,1)\)，递减区间为 \((1,+\infty)\). 

（2）因为
\(b\ln a-a\ln b=a-b\)，
故
\(b\left( \ln a+1 \right)=a\left( \ln b+1 \right)\)，
即
\(\frac{\ln a+1}{a}=\frac{\ln b+1}{b}\).

于是
\(f\left( \frac{1}{a} \right)=f\left( \frac{1}{b} \right)\).

设
\(\frac{1}{a}=x_1\)，\(\frac{1}{b}=x_2\).
由（1）可知，不妨设
\(0<x_1<1\)，\(x_2>1\).

因为当 \(x\in(0,1)\) 时，
\(f(x)=x\left( 1-\ln x \right)>0\)，
当 \(x\in(\e,+\infty)\) 时，
\(f(x)=x\left( 1-\ln x \right)<0\)，
故
\(1<x_2<\e\).

先证
\(x_1+x_2>2\).

若 \(x_2\ge 2\)，则由 \(0<x_1<1\) 可知 \(x_1>0\)，从而 \(x_1+x_2>2\). 

若 \(x_2<2\)，则要证 \(x_1+x_2>2\)，只需证
\(x_1>2-x_2\).
而
\(0<2-x_2<1\)，
故只需证
\(f(x_1)>f(2-x_2)\).

又
\(f(x_1)=f(x_2)\)，
所以只需证
\(f(x_2)>f(2-x_2)\).

设
\(g(x)=f(x)-f(2-x)\)，\(1<x<2\).
则
\(g'(x)=f'(x)+f'(2-x)=-\ln x-\ln(2-x)=-\ln\left[ x(2-x) \right]\).

因为 \(1<x<2\)，故
\(0<x(2-x)<1\)，
所以
\(g'(x)>0\).
故 \(g(x)\) 在 \((1,2)\) 上为增函数，所以
\(g(x)>g(1)=0\).

因此
\(f(x_2)>f(2-x_2)\)，
从而
\(x_1+x_2>2\).

再证
\(x_1+x_2<\e\).

设
\(x_2=tx_1\)，\(t>1\).
结合
\(\frac{\ln a+1}{a}=\frac{\ln b+1}{b}\)
与
\(\frac{1}{a}=x_1\)，\(\frac{1}{b}=x_2\)，
可得
\(x_1\left( 1-\ln x_1 \right)=x_2\left( 1-\ln x_2 \right)\).

即
\(1-\ln x_1=t\left( 1-\ln t-\ln x_1 \right)\)，
解得
\(\ln x_1=\frac{t-1-t\ln t}{t-1}\).

要证 \(x_1+x_2<\e\)，即证
\(\left( t+1 \right)x_1<\e\)，
即证
\(\ln(t+1)+\ln x_1<1\).

代入 \(\ln x_1\) 的表达式，等价于证
\(\left( t-1 \right)\ln(t+1)-t\ln t<0\).

令
\(S(t)=\left( t-1 \right)\ln(t+1)-t\ln t\)，\(t>1\).
则 \(S'(t)=\ln(t+1)+\frac{t-1}{t+1}-1-\ln t=\ln\left( 1+\frac{1}{t} \right)-\frac{2}{t+1}\).

先证明不等式
\(\ln(1+x)\le x\).
设
\(u(x)=\ln(1+x)-x\)，
则
\(u'(x)=\frac{1}{x+1}-1=-\frac{x}{x+1}\).
当 \(-1<x<0\) 时，\(u'(x)>0\)；当 \(x>0\) 时，\(u'(x)<0\)，故
\(u(x)\le u(0)=0\)，
即
\(\ln(1+x)\le x\).

于是当 \(t>1\) 时，
\(\ln\left( 1+\frac{1}{t} \right)\le \frac{1}{t}<\frac{2}{t+1}\)，
故
\(S'(t)<0\).
所以 \(S(t)\) 在 \((1,+\infty)\) 上为减函数，从而
\(S(t)<S(1)=0\).

因此
\(x_1+x_2<\e\).

综上所述，
\(2<\frac{1}{a}+\frac{1}{b}<\e\).
\end{solution}
